\structuralpart{ПРИЛОЖЕНИЕ А}
\begin{center}Основные исходные тексты\end{center}
\renewcommand{\thelstlisting}{А.\arabic{lstlisting}}
\setcounter{lstlisting}{0}
В листинге~\ref{lst:api} «Интерфейс функций сегментации» приведены объявления типов и функций. В листинге~\ref{lst:implementation} «Реализация метода Оцу и управления потоками» приведён исходный код вычислительного модуля, использованного при тестировании и измерениях.
\lstinputlisting[caption={Интерфейс функций сегментации},label={lst:api}]{src/segmentation.h}
\clearpage
\lstinputlisting[caption={Реализация метода Оцу и управления потоками},label={lst:implementation}]{src/segmentation.c}

\structuralpart{ПРИЛОЖЕНИЕ Б}
\begin{center}Полные результаты измерений\end{center}
\renewcommand{\thetable}{Б.\arabic{table}}\setcounter{table}{0}
В таблице~\ref{tab:all} представлены сводные результаты всех конфигураций. В столбце $q$ обозначением «Послед.» отмечена самостоятельная последовательная реализация. Для каждой строки учтено 100 измерений, время приведено в миллисекундах. Минимум и максимум характеризуют отдельные вызовы, а не доверительный интервал среднего.
\begin{longtable}{|p{23mm}|p{29mm}|p{26mm}|p{26mm}|p{26mm}|}
\caption{Полные результаты измерений}\label{tab:all}\\\hline
$q$ & Среднее & Станд. откл. & Минимум & Максимум\\\hhline{|=:=:=:=:=|}
\endfirsthead
\multicolumn{5}{l}{Продолжение таблицы~\thetable}\\\hline
$q$ & Среднее & Станд. откл. & Минимум & Максимум\\\hhline{|=:=:=:=:=|}
\endhead
\multicolumn{5}{|c|}{Изображение $256 \times 256$}\\\hline
Послед. & 0,186 & 0,020 & 0,174 & 0,251 \\\hline
1 & 0,401 & 0,085 & 0,327 & 0,826 \\\hline
2 & 0,342 & 0,071 & 0,267 & 0,628 \\\hline
4 & 0,404 & 0,095 & 0,322 & 0,887 \\\hline
8 & 0,586 & 0,081 & 0,484 & 0,897 \\\hline
16 & 1,015 & 0,217 & 0,816 & 2,488 \\\hline
32 & 1,995 & 0,570 & 1,586 & 5,907 \\\hline
\multicolumn{5}{|c|}{Изображение $512 \times 512$}\\\hline
Послед. & 0,701 & 0,046 & 0,670 & 0,869 \\\hline
1 & 1,075 & 0,122 & 1,000 & 2,158 \\\hline
2 & 0,724 & 0,167 & 0,593 & 1,702 \\\hline
4 & 0,716 & 0,187 & 0,620 & 2,181 \\\hline
8 & 0,832 & 0,115 & 0,742 & 1,486 \\\hline
16 & 1,166 & 0,187 & 1,018 & 1,990 \\\hline
32 & 1,982 & 0,296 & 1,730 & 2,963 \\\hline
\multicolumn{5}{|c|}{Изображение $1024 \times 1024$}\\\hline
Послед. & 3,538 & 1,904 & 2,654 & 9,319 \\\hline
1 & 4,900 & 2,750 & 3,615 & 13,763 \\\hline
2 & 3,242 & 1,841 & 1,897 & 9,761 \\\hline
4 & 2,638 & 1,473 & 1,874 & 9,245 \\\hline
8 & 2,793 & 1,544 & 1,983 & 8,101 \\\hline
16 & 3,123 & 1,766 & 2,242 & 13,246 \\\hline
32 & 4,011 & 2,020 & 2,786 & 11,782 \\\hline
\multicolumn{5}{|c|}{Изображение $2048 \times 2048$}\\\hline
Послед. & 34,857 & 6,053 & 16,462 & 45,279 \\\hline
1 & 46,870 & 8,813 & 22,562 & 81,897 \\\hline
2 & 30,017 & 7,893 & 12,983 & 60,721 \\\hline
4 & 21,900 & 4,963 & 10,222 & 47,354 \\\hline
8 & 20,732 & 4,466 & 9,916 & 44,206 \\\hline
16 & 20,913 & 4,471 & 9,918 & 44,469 \\\hline
32 & 22,526 & 7,737 & 11,221 & 85,670 \\\hline
\multicolumn{5}{|c|}{Изображение $4096 \times 4096$}\\\hline
Послед. & 132,034 & 21,368 & 70,160 & 182,776 \\\hline
1 & 173,095 & 25,124 & 100,864 & 259,892 \\\hline
2 & 102,399 & 19,625 & 60,640 & 150,750 \\\hline
4 & 76,934 & 12,599 & 48,247 & 130,949 \\\hline
8 & 73,955 & 10,001 & 49,148 & 91,837 \\\hline
16 & 74,901 & 8,838 & 45,157 & 90,180 \\\hline
32 & 76,335 & 8,286 & 57,570 & 103,123 \\\hline
\end{longtable}
